\chapter{Conclusion}

Le projet ITI Aventure nous a permis de mettre en pratique plusieurs compétences acquises durant l’année, en particulier la persistance des données avec JDBC, la création d’un interpréteur via JavaCC, et l’utilisation de design patterns pour améliorer la structure du code.

Nous avons conçu une base relationnelle adaptée au jeu, développé un système de lecture et d’écriture des données en Java, implémenté un langage de description fonctionnel pour créer un monde, et généralisé notre code pour permettre des extensions comme ITISpaceOpera.

Au-delà des aspects techniques, ce projet à développer notre sens de la communication en équipe, à résoudre des problèmes concrets de développement logiciel, et à concevoir une architecture claire et évolutive. Il a constitué une expérience plutôt riche, d'une part par le fait qu'il s'agit du projet informatique le plus concret de l'année et d'autre part car le temps très limité à notre disposition nous contraignait à être prêt pour la date de rendu malgré notamment les contraintes de révision des examens.